% !TeX root = ../main.tex
\begin{frame}{课程评价——王思霖}
  \begin{SlideCanvas}
    \node[anchor=north west,inner sep=0pt,text width=850pt,align=left,
      text=BICNavy] at (55,80) {\begin{minipage}{850pt}
      \usebeamerfont{BIC detail}\fontsize{18}{22}\selectfont\raggedright
      \setlength{\parindent}{0pt}\setlength{\parskip}{3pt}
      \setlength{\columnsep}{26pt}\setlength{\multicolsep}{0pt}
      \clubpenalty=10000\widowpenalty=10000
\begin{multicols}{2}\raggedcolumns
\par\addvspace{6pt}{\usebeamerfont{BIC point}一、课程评价\par}\nobreak
本课程邀请医学、医疗器械、集成电路、机械工程等多领域专家开展授课，内容涵盖医学影像、内镜诊疗、医疗传感、智能手术机器人等前沿研究方向。课程紧密结合临床需求与工程落地场景，内容体系详实、案例贴合实际，有效拓宽了研究生的医工交叉科研视野。\par
\par\addvspace{6pt}{\usebeamerfont{BIC point}二、学习收获\par}\nobreak
作为机械方向的研究生，通过本课程学习，我初步建立起医工交叉的研究思维，拓展了跨学科科研视野。课程中的各类临床案例，让我了解了内镜微创、神经外科诊疗过程中存在的实际难题。这类高精度、精细化的医疗操作，对设备结构设计、运动控制与操作稳定性都有着严苛要求。\par
同时我了解到，当代医疗装备正逐步从传统刚性结构向刚柔耦合机器人、智能感知与人机交互等方向升级发展。依托机械结构优化、精密传动设计与智能控制技术，智能手术机器人能够突破现有医疗设备的技术瓶颈，进一步提升临床诊疗的精准度与安全性。\par
此外，我也认识到当前医工交叉发展存在的难点：医学与工程学科的专业体系差异较大，临床需求难以快速、精准转化为工程设计方案；新型医疗设备研发周期长、落地成本高。这也让我明白，医工交叉研究必须以临床实际需求为核心，立足实际问题开展技术创新。\par
\par\addvspace{6pt}{\usebeamerfont{BIC point}三、意见与建议\par}\nobreak
1.~增设设备演示教学环节。通过超声探头、内镜器械、医疗机器人样机的实物展示或仿真演示，帮助学生直观理解医疗设备的结构原理与工作机制；\par
2.~提前推送课程资料。针对医学、医疗器械等跨学科内容，课前发放基础导读材料，有效降低跨专业学习门槛，进一步提升课堂学习效率；\par
3.~增设跨专业学生交流环节。组织集成电路、机械等不同专业学生开展小组讨论，结合各自学科视角交流讨论，促进多学科思维碰撞与交叉协作。\par
\end{multicols}
    \end{minipage}};
  \end{SlideCanvas}
  \note{原始材料为用户提供的四份课程反馈之一，正文逐字保留，仅调整标题、段落和编号排版。王思霖《医工交叉研究前沿进展》课程反馈与建议。一、课程评价\par 本课程邀请医学、医疗器械、集成电路、机械工程等多领域专家开展授课，内容涵盖医学影像、内镜诊疗、医疗传感、智能手术机器人等前沿研究方向。课程紧密结合临床需求与工程落地场景，内容体系详实、案例贴合实际，有效拓宽了研究生的医工交叉科研视野。\par 二、学习收获\par 作为机械方向的研究生，通过本课程学习，我初步建立起医工交叉的研究思维，拓展了跨学科科研视野。课程中的各类临床案例，让我了解了内镜微创、神经外科诊疗过程中存在的实际难题。这类高精度、精细化的医疗操作，对设备结构设计、运动控制与操作稳定性都有着严苛要求。\par 同时我了解到，当代医疗装备正逐步从传统刚性结构向刚柔耦合机器人、智能感知与人机交互等方向升级发展。依托机械结构优化、精密传动设计与智能控制技术，智能手术机器人能够突破现有医疗设备的技术瓶颈，进一步提升临床诊疗的精准度与安全性。\par 此外，我也认识到当前医工交叉发展存在的难点：医学与工程学科的专业体系差异较大，临床需求难以快速、精准转化为工程设计方案；新型医疗设备研发周期长、落地成本高。这也让我明白，医工交叉研究必须以临床实际需求为核心，立足实际问题开展技术创新。\par 三、意见与建议\par 1.~增设设备演示教学环节。通过超声探头、内镜器械、医疗机器人样机的实物展示或仿真演示，帮助学生直观理解医疗设备的结构原理与工作机制；\par 2.~提前推送课程资料。针对医学、医疗器械等跨学科内容，课前发放基础导读材料，有效降低跨专业学习门槛，进一步提升课堂学习效率；\par 3.~增设跨专业学生交流环节。组织集成电路、机械等不同专业学生开展小组讨论，结合各自学科视角交流讨论，促进多学科思维碰撞与交叉协作。}
\end{frame}
